\section{The ordinary XY model: mixed Gaussian histories}
\label{xyh:section}

This section constructs the expansion used for the ordinary cosine interaction.
The sizes constructed below are sizes of specified integral representations.
Their defining estimate concerns an evaluated history, and does not assert a
pointwise absolute bound at every value of an unintegrated complex field.
This distinction is necessary because the compensating gradient interaction
has positive sign.

The aim is the composable activity calculus of
Proposition~\ref{xyh:calculus}, used by the recurrence in
Section~\ref{xyf:section}.  The branch identity supplies the
microscopic factors, and the local operations assemble them into the
selected histories of Definition~\ref{xyh:history-definition}.
The support and atom estimates, the evaluated-history bound of
Lemma~\ref{xyh:endpoint-estimate}, and the separate charge payments
of Lemma~\ref{xyh:budget} together give the small costs and
single-input gains in that calculus.

The Fourier and vortex representations belong to the duality
framework of \citet{JoseKadanoffKirkpatrickNelson1977}.
The Gaussian-smoothed branch decomposition and the evaluated-history
bounds below supply the additional estimates for the cosine interaction.

\subsection{An exact branch representation}

Use lattice units, fuse the boundary of the spin square to one vertex,
and delete the pinned perimeter edges, as in Lemma~\ref{vi:duality}.
Each deleted edge has weight $e^{b\cos0}=e^b$, or one after the
normalization by $e^{-b}$ used below, so this deletes only a
source-independent factor.
Let $d$ and $\partial$ be the primal and dual incidence operators, and let
$\mathcal R$ be the signed crossing identification, with the orientations of
Lemma~\ref{vi:duality}. Inner products count every edge once. Put
\[
 G=(\partial^*\partial)^{-1},\qquad
 P_L=\partial G\partial^*,\qquad P_T=I-P_L,
 \qquad \mathfrak p=2\pi\sqrt b,
\]
where the inverse is on dual scalars of mean zero. Thus
$P_T=\mathcal R^{-1}dG_n d^*(\mathcal R^{-1})^*$.
Fix a sufficiently small number $\upsilon>0$, independent of $b$, and define
\[
 h_b(x)=\int_{-\pi\sqrt b}^{\pi\sqrt b}
 \frac{e^{-(x-t)^2/(2\upsilon^2)}}{\sqrt{2\pi}\upsilon}\,dt,
 \qquad
 f_b(z)=h_b(z)e^{z^2/2+b(\cos(z/\sqrt b)-1)},\qquad F_b=f_b-1.
\]
Both functions in the last display are entire. The intervals translated by
$\mathfrak p\Z$ partition the real line, up to their endpoints, so
$\sum_{m\in\Z}h_b(x+\mathfrak p m)=1$. Consequently
\begin{equation}\label{xyh:periodization}
 \sum_{m\in\Z} e^{-(x+\mathfrak p m)^2/2}
 f_b(x+\mathfrak p m)=e^{b(\cos(x/\sqrt b)-1)}.
\end{equation}
In particular smoothing the branch cutoff has made no change to the spin
law.

\begin{lemma}[Bounds for the branch factor]\label{xyh:branch-bounds}
There are $b_1,C<\infty$ and $\kappa_0<1/2$ such that, for $b\ge b_1$,
\begin{equation}\label{xyh:primitive-bound}
 |F_b(x+iy)|\le \frac Cb
  \exp\{\kappa_0 x^2+C y^2 e^{|y|/\sqrt b}\},
 \qquad (\log f_b)''(x)\ge-C\quad(x,y\in\R).
\end{equation}
\end{lemma}
\begin{proof}
On the central interval, $1-\cos(x/\sqrt b)\ge c x^2/b$.
On a fixed slightly larger interval the same inequality holds with a smaller
$c>0$. Outside that interval the Gaussian tail of $h_b$ gives the bound
$h_b(x)e^{b(\cos(x/\sqrt b)-1)}\le C e^{-c x^2}$, after choosing
$\upsilon$ small. This proves the bound globally on the real line. Directly
from the integral defining $h_b$,
$|h_b(x+iy)|\le e^{y^2/(2\upsilon^2)}h_b(x)$; also
$b|\cos(x/\sqrt b)|(\cosh(y/\sqrt b)-1)
\le C y^2 e^{|y|/\sqrt b}$. Hence
$|f_b(x+iy)|\le C e^{(1/2-c)x^2+C y^2 e^{|y|/\sqrt b}}$.

To obtain the factor $b^{-1}$, on $|x|\le\sqrt b$ write $f_b-1$ as
$(h_b-1)e^{z^2/2+b(\cos(z/\sqrt b)-1)}$
plus $e^{z^2/2+b(\cos(z/\sqrt b)-1)}-1$.
The first term uses the exponentially small cutoff error, with the same
imaginary bound. Taylor's formula with integral remainder bounds the
exponent in the second by
$Cb^{-1}(|x|+|y|)^4e^{|y|/\sqrt b}$.
The inequality $|e^u-1|\le |u|\max(1,e^{\Re u})$ and the previously
proved real Gaussian margin absorb the $x$ polynomial into a fixed
increase of its square coefficient; the $y$ polynomial is absorbed into
$C e^{C y^2e^{|y|/\sqrt b}}$. On $|x|>\sqrt b$, the global bound
and $b\le C e^{c x^2/2}$ give the same result. These increases leave
$\kappa_0<1/2$. Finally differentiation under the positive real integral
shows
\[
 (\log h_b)''(x)=-\upsilon^{-2}
 +\upsilon^{-4}\Var_x(t)\ge-\upsilon^{-2}.
\]
Adding the second derivative $1-\cos(x/\sqrt b)\ge0$ of the other
exponent proves the Hessian bound.
\end{proof}

\begin{proposition}[Unwrapping and sources]\label{xyh:branch}
Up to a positive source-independent constant, the XY partition function is
\begin{equation}\label{xyh:positive-ensemble}
 \sum_{\substack{q\in\Z^{\mathrm{dual}}\\\sum q=0}}
 e^{-\|V_q\|^2/2}\,
 \E_{X\sim N(0,P_T)}\prod_e f_b((X+V_q)_e),
 \qquad V_q=\mathfrak p\partial Gq.
\end{equation}
For an integral source $l$ on primal vertices, choose an integral current
$dG_n l+\mathcal R\partial a$ of divergence $l$, and set
\[
 \vartheta=a/\sqrt b,\qquad
 \eta=-\mathcal R^{-1}dG_n l/\sqrt b.
\]
The numerator for $\E e^{i(l,\theta)}$ is the same sum with insertion
\begin{equation}\label{xyh:physical-insertion}
 e^{-i(X,\eta)+i\mathfrak p(q,\vartheta)}.
\end{equation}
In particular $\partial^*\eta=0$ and
$\|\eta\|^2=(l,G_nl)/b$.
\end{proposition}
\begin{proof}
Apply Equation~\eqref{xyh:periodization} on each edge, with
$x=\sqrt b\,d\theta$. For a branch integer $m$, put
$w=\mathcal R^{-1}\sqrt b(d\theta+2\pi m)$.
Its divergence is $\partial^*w=\mathfrak p q$ for an integral mean-zero
$q$. The incidence matrix of a connected graph maps integral edges onto
integral mean-zero vertex arrays: a rooted spanning tree gives a solution
by assigning to each parent edge the sum of charges below it. Thus every
such $q$ occurs. Two integral branch solutions with the same $q$ differ
by an integral primal gradient after crossing. To verify this last claim,
integrate their increments along primal paths; zero sums around dual
plaquettes make the result path independent, and fixing its value at the
fused vertex makes it integral and unique. Summing over these gradients
unwraps the angular fundamental domain to the complete transverse real
space. Its Jacobian is constant, independent of $q$ and of the source.
Orthogonal decomposition gives $w=X+V_q$, with $X\in\operatorname{ran}P_T$;
its Gaussian square splits into $\|X\|^2+\|V_q\|^2$.

In the angular insertion the continuous lift gives $-i(X,\eta)$.
Changing that lift by the gradient part of $2\pi m$ contributes
$-2\pi(dG_nl,m)=2\pi(q,a)$ modulo $2\pi$, because the selected total
current is integral. This is exactly the second factor in
Equation~\eqref{xyh:physical-insertion}. Orthogonality and the definition
of $G_n$ give the last two assertions. All sums converge on a finite graph
by the real Gaussian branch bound in Lemma~\ref{xyh:branch-bounds}.
\end{proof}

\subsection{Compensation and the Gaussian prescription}

Write $p=\partial\psi$ and use the parameter notation
\[
 \mu=1+s_{11},\qquad d_c=1+s_{12},\qquad s=s_{22},\qquad
 U(w,p)=\tfrac12\sum_e(s_{11}w_e^2+2i s_{12}w_ep_e+s p_e^2).
\]
The compensation parameters $(s_{11},s_{12},s_{22})$ lie in a fixed
sufficiently small neighborhood of zero and satisfy
\begin{equation}\label{xyh:permitted-parameters}
 s\ge d_*s_{12}^2,\qquad d_*>0\text{ fixed}.
\end{equation}
Before compensation the inverse covariance is $\mu$ on transverse $w$
and
\begin{equation}\label{xyh:inverse-pair}
 \begin{pmatrix}\mu&i d_c\\i d_c&s\end{pmatrix}
\end{equation}
on longitudinal pairs $(w,p)$. Scalars $\psi$ also have a uniform common
constant of period $b^{-1/2}$. The vertex factors are
$\sum_{q_v\in\Z}e^{i\mathfrak p q_v\psi_v}$; this common constant
enforces $\sum q_v=0$.

\begin{proposition}[Exact compensation and spin offsets]
\label{xyh:compensation}
The normalized integral with inverse covariance
Equation~\eqref{xyh:inverse-pair}, vertex factors, edge factors $f_b(w)$,
and factor $e^{U(w,p)}$ equals
Equation~\eqref{xyh:positive-ensemble} up to one source-independent
determinant normalization. If in the interaction one replaces
\begin{equation}\label{xyh:offsets}
 p\mapsto p+\eta,\quad \psi\mapsto\psi+\vartheta,
 \quad w\mapsto w-i d_c\eta/\mu,
\end{equation}
including this last replacement in $f_b$, the integral is the physical
inserted partition function multiplied by
\begin{equation}\label{xyh:spin-prefactor}
 \exp\{(s+d_c^2/\mu)\|\eta\|^2/2\}.
\end{equation}
These identities have the finite-dimensional limiting interpretation
specified in the proof, including at $s=0$.
\end{proposition}
\begin{proof}
At first take $s>0$, finitely many Fourier indices, and multiply every
compensating quadratic by a number $\lambda<1$. Integration then uses the
ordinary complex Gaussian density whose real precision on $w,p$ is
positive. Normalize its Gaussian integral by its determinant, continuously
from real positive diagonal precision. The mixed off-diagonal entries are
purely imaginary, so this prescription is unambiguous. On restoring the
full compensation its formal exponent becomes
$-\|w\|^2/2-i(w,p)$. Integrating a fixed neutral Fourier charge by Gaussian
completion gives real mean $V_q$ and covariance $P_T$ for $w$; the
remaining charge factor is $e^{-\|V_q\|^2/2}$. The matrix calculation is
also obtained by first integrating $p$ at $\lambda<1$ and letting
$\lambda\uparrow1$; the covariance tends to the stated positive
semidefinite matrix. The determinant ratio is independent of $q$.

For the offset identity put $w'=w-i d_c\eta/\mu$ in the full compensated
exponent. Expansion gives exactly
\[
 -\tfrac12\|w'\|^2-i(w',p)-i(w',\eta)
       +\tfrac12(s+d_c^2/\mu)\|\eta\|^2.
\]
Every additional linear term involving $p$ vanishes since $(p,\eta)=0$.
The vertex shift supplies $e^{i\mathfrak p(q,\vartheta)}$.
The contour translation in $w$ is allowed at every fixed regularization
by the strict real Gaussian margin of
Equation~\eqref{xyh:primitive-bound}; for an entire integrand the vertical
sides of a translating rectangle have integrals tending to zero.

For completeness the uniform domination needed here, and in subsequent
local calculations, is proved in Lemma~\ref{xyh:endpoint-estimate} and
Proposition~\ref{xyh:calculus}: at fixed depth it retains
$e^{-c b\sum q_v^2}$ after any fixed deterministic contour shifts and
permits $\lambda\uparrow1$ and Fourier truncation removal. The use of that
estimate is noncircular: its proof starts with the same strict, finite
regularizations and only uses their finite Gaussian completion identities.
It proves domination before asserting the present limiting identities.
At $s=0$, Equation~\eqref{xyh:permitted-parameters} forces $s_{12}=0$;
take a limit from the permitted region. Source-independent determinants
cancel in every partition ratio. Thus no positive joint probability law
for $(w,\psi)$ is being assumed at a degenerate endpoint.
\end{proof}

\subsection{Kernels and normalized smooth tests}

The finite grids, their reflections, and block geometry are those in
Lemmas~\ref{rg:kernels}, \ref{rg:green}, and \ref{rg:geometry}. We modify
only the low-pass polynomial in order to allow increasing differentiability.
Let $p_0=32$, and, for dyadic $d$, put
\[
 P_d(\cos t)=\left(\frac{\sin(dt/2)}{d\sin(t/2)}\right)^{2p_0},
 \qquad Q_r=\prod_{\substack{2\le d\le r\\d\text{ dyadic}}}P_d,
 \qquad Q_1=1.
\]
The arguments of these polynomials are $1-\lambda/16$, where $\lambda$
is the graph-Laplacian eigenvalue. After the first piece $\gamma I$,
$\gamma=1/16$, the binary scalar pieces have multipliers
\begin{equation}\label{xyh:binary-multiplier}
 \widehat\Gamma_r(\lambda)
  =(1-\gamma\lambda)\frac{Q_r-Q_{2r}}{\lambda}.
\end{equation}
Use the same prescription on dual free and primal Dirichlet scalars,
writing $\Gamma_r^{\mathrm{pr}}$ for the latter. The terminal positive
remainder is split, for estimates, into binary pieces through a size
comparable to the domain and a final remainder. On the auxiliary torus
include the two independent harmonic transverse directions in this last
remainder. Harmless scalar constant variances may also be added there.

\begin{lemma}[High-order kernel estimates]\label{xyh:kernels}
There are constants $C,C_0,c>0$, independent of the shell ratio $L$, such
that, at binary scale $r$ and for every difference order $j$,
\begin{align}
 |\nabla^j\Gamma_r(x,y)|&\le C e^{C_0j^2}r^{-j},
 &\|\nabla^j\Gamma_r^{1/2}\|_{2\to2}
  &\le C e^{C_0j^2}r^{1-j}\quad(j\ge1).
 \label{xyh:high-derivatives}
\end{align}
The kernel has range $Cr+O(j)$. These bounds hold for the reflected kernels
and for the last remainder on nonconstant modes, with $r$ comparable to
the domain size. A shell from $m$ to $Lm$ has diagonal at least $c\log L$
on the free grid, plane, or torus, and its positive-order differences on
a region of diameter $O(m)$ obey the bounds obtained by summing
Equation~\eqref{xyh:high-derivatives}.

Set $D_r=\partial\Gamma_r^{1/2}$ and
$B_r=\mathcal R^{-1}d(\Gamma_r^{\mathrm{pr}})^{1/2}$. Including first
and last pieces and the stipulated harmonic directions,
\begin{equation}\label{xyh:white-resolution}
 E_r=B_rB_r^*+D_rD_r^*\succeq0,\quad \sum_rE_r=I,
 \quad E_{\ge S}:=I-\sum_{r<S}E_r
 \text{ has range }CS.
\end{equation}
\end{lemma}
\begin{proof}
Each $P_d$ is a nonnegative polynomial at most one on the spectrum, of
degree $p_0(d-1)$. Thus $Q_r-Q_{2r}=Q_r(1-P_{2r})\ge0$ and its zero
at $\lambda=0$ makes Equation~\eqref{xyh:binary-multiplier} a polynomial
of degree $O(r)$. This proves positivity and range. On the Fourier square
$\lambda\asymp|\xi|^2$, and the product over dyadic $d$ gives
\[
 |\widehat\Gamma_r(\xi)|\le
 C\min\{r^2,\ |\xi|^{-2}
       e^{-c(\log^+(r|\xi|))^2}\}.
\]
Indeed for $d|\xi|\ge1$ the factor is bounded by
$C(d|\xi|)^{-2p_0}$, whose logarithms sum to a negative quadratic in
$\log(r|\xi|)$; for $r|\xi|\le1$ use
$1-P_{2r}\le Cr^2|\xi|^2$. Multiplication by $|\xi|^j$ and integration
in two dimensions reduces the pointwise estimate to
$r^{-j}\int_1^\infty t^{j-1}e^{-c(\log t)^2}\,dt$, bounded by
$C e^{C_0j^2}r^{-j}$ by completing the square. Taking a supremum instead
of an integral proves the square-root operator estimate. Discrete Fourier
sums give the same bounds on finite periodic extensions. Reflection has
only a fixed number of images within range at scales below the domain;
the terminal Fourier sum gives the remaining cases. For
$j\le8+\lceil\log(2+r)\rceil$, the stencil extension is still $O(r)$.

On a small annular frequency band $a/r\le |\xi|\le4a/r$, the earlier
product is bounded below since its relative errors sum as
$\sum_{d\le r}d^2|\xi|^2\le C a^2$; the last difference is bounded
below by $c_a r^2$. Each binary piece consequently contributes a positive
constant to its diagonal. For reflected cosine eigenfunctions, pair a
frequency index $u$ with $2u$ inside a slightly wider band: at least one
of their squared cosines is bounded below at every site. Apply this in
both coordinates. The terminal clearance ensures enough frequencies per
band. Summing the binary scales proves the diagonal assertion.

The scalar pieces telescope to $G$ and $G_n$, respectively. Applying
gradients yields $P_L+P_T=I$; on a torus the harmonic addition supplies
the missing two directions. Finally $E_{\ge S}$ is the identity minus
a finite sum of polynomial kernels of total range $CS$.
\end{proof}

The same binary summation proves the high-order version of the Green
bounds used for local lifts:
\begin{equation}\label{xyh:green-high}
 |\nabla_x^jG_n(x,y)|\le C e^{C_0j^2}
             (1+\dist(x,y))^{-j},\qquad j\ge1.
\end{equation}
Indeed pieces of size smaller than a fixed multiple of the distance
vanish, except for the difference-stencil enlargement; the remaining
geometric sum of $r^{-j}$ gives the bound. When the stencil length is
comparable to the distance, its extra factor $j^j$ is absorbed by
$e^{C_0j^2}$. The box-scale last piece has the same bound. Differentiating
$\eta$ and integrating it on a source-free local region yields the
corresponding test bounds for local lifts with these high-order
constants. Reflections use the oriented parities of the primal and dual
kernels. This extends the fixed-order bounds of
Lemma~\ref{vi:source-lifts} to every order needed here.

At scale $m$ let $j(m)=8+\lceil\log(2+m)\rceil$. A normalized scalar
test on a specified region and fixed collar satisfies
\begin{equation}\label{xyh:test-class}
 |m^j\nabla^j u|\le e^{K_0j^2}\quad(0\le j\le j(m));
\end{equation}
a normalized edge test satisfies the same inequalities for $m w$.
Here $K_0$ is fixed larger than all high-order kernel constants. Functions
are defined on the entire extended grid, but only the stated local bounds
are imposed; derivatives at walls use reflection. A scalar atom born at
$m$ is a linear functional supported in a fixed multiple of an $m$-box,
annihilating constants, and bounded by one on these scalar tests. An edge
atom uses the edge test convention, without constant annihilation.
Multiplication by any fixed constant merely changes the universal atom
bound. For a real charge array $k$ on $X$, write
\[
 h_m(k;X)=\sup_{u\text{ satisfying }\eqref{xyh:test-class}}
                |(k,u)|,
\]
using a fixed expanded neighborhood of $X$. Coarse tests restrict to fine
tests. If $X$ is small, subtraction of a constant gives
\begin{equation}\label{xyh:test-coarsening}
 h_{Lm}(k;\overline X)\le h_m(k;X),\qquad
 h_{Lm}(k;\overline X)\le |\textstyle\sum k|+(C/L)h_m(k;X).
\end{equation}
The second inequality follows by joining each site to an anchor with a
path of length $O(m)$, and applying the positive-order test inequalities.

The next four lemmas prepare the estimates for a sequence of local
integrations and Taylor subtractions.  One term selected from this
expansion records which microscopic factors were used, which components
merged, and which Taylor directions were chosen.  Its precise record,
the selected history of Definition~\ref{xyh:history-definition}, is a
rooted forest: microscopic factors are leaves, and operations join or
end their field dependence.  A Taylor subtraction on one component
introduces atoms; evaluating a scalar coefficient ends the old field
dependence.  The forest rules will be specified after the analytic
estimates, then used to check their overlap and scale-weight hypotheses
separately on each coupled live root.

First we separate a coarse Taylor argument into fixed normalized fine
directions and linear coarse coefficients.  Lemma~\ref{xyh:fourier}
gives small absolute coefficient sums for this separation.  Those sums
will supply the powers of $L^{-1}$ in the local Taylor gains; the atoms
remain as field-dependent factors to be controlled by the subsequent
lemmas.

\begin{lemma}[Smooth directional decomposition]\label{xyh:fourier}
Let $S=Lm$, and restrict a coarse field $z=(w,\psi)$ to the
prescribed neighborhood of a small fine activity, with its fine
collar. Modulo a scalar constant, this restriction is a finite sum
of fixed normalized fine test directions multiplied by bounded
$S$-atom evaluations, with total absolute numerical coefficient
weight at most $C/L$. The bulk error after subtracting a constant
oriented edge field and an affine scalar field has such a
representation with total weight at most $C/L^2$.
The constants hold for all sufficiently large $L$ and for the
specified reflected and unwrapped periodic regions.
\end{lemma}
\begin{proof}
Use the test convention in Equation~\eqref{xyh:test-class}, and put
\[
 W(t)=\sup_{j\in\mathbb N_0}t^j e^{-K_0j^2}.
\]
Choose a periodic lattice box of side $T\asymp S$ containing the
fine region and its collar. A cutoff $\chi$ equals one there, is
supported in the prescribed coarse neighborhood, and satisfies
\[
 |\nabla^k\chi|\le C^{k+1}(k!)^3S^{-k}.
\]
Such a cutoff is obtained by sampling a fixed compact Gevrey bump
after rescaling by $S$. For scalar fields expand
$f=\chi(\psi-\psi(v_0))$, and for edge fields expand $f=S\chi w$.
Write $c_\ell(f)$ for normalized discrete Fourier coefficients and
$r_\ell=1+|\ell|$, using representatives with $|\ell_i|\le T/2$.

The cutoff preserves the derivative allowance at the cost of an
exponential factor in the order. Indeed, a Leibniz term with $k$
derivatives on the cutoff has, relative to $e^{K_0j^2}$, the factor
\[
 (k!)^3 e^{K_0(j-k)^2-K_0j^2}
 \le (k!)^3e^{-K_0k^2}\le C_{K_0}.
\]
Summing the Leibniz terms, including harmless lattice shifts, gives
$S^j|\nabla^jf|\le C^{j+1}e^{K_0j^2}$ for $0\le j\le j(S)$
on normalized coarse tests. Summation by parts in a coordinate
with maximal $|\ell_i|$ gives
\[
 |c_\ell(f)|\le C C_1^j r_\ell^{-j}e^{K_0j^2},
 \qquad 0\le j\le j(S).
\]
For sufficiently small fixed $c>0$,
\[
 A_\ell(z)=W(cr_\ell)c_\ell(f)
\]
is therefore a bounded coarse atom. Scalar coefficients annihilate
constants because of the anchor subtraction. The finite derivative
budget suffices: a maximizing index for $W(t)$ can be chosen at most
$J$ whenever $\log t\le K_0(2J+1)$; this holds for $t=cr_\ell$,
$|\ell|\le CT$, and $J=j(S)$.

Let $e_\ell$ denote a Fourier mode. On the fine region the scalar
mode $e_\ell-e_\ell(v_0)$ has fine test norm at most
\[
 C L^{-1}r_\ell W(Cr_\ell/L).
\]
Subtracting its discrete affine interpolant gives the bound
\[
 C L^{-2}r_\ell^2W(Cr_\ell/L).
\]
At orders zero and one this is discrete Taylor's formula; at orders
at least two the two frequency factors are already present.
For edges the Fourier mode carries the factor $S^{-1}$, and the
corresponding polynomial powers of $r_\ell$ are zero and one.
The cutoff equals one at every site needed for the anchor and
its difference stencils, so these are exact identities for the
specified restrictions and affine errors.

For every fixed positive integer $N$, shifting the integer in the
supremum defining $W$ gives
\[
 W(t)\ge e^{-K_0N^2}t^N W(te^{-2K_0N}).
\]
If $L\ge(C/c)e^{2K_0N}$, then
\[
 \frac{W(Cr/L)}{W(cr)}\le e^{K_0N^2}(cr)^{-N}.
\]
Choose $N>P+2$, where $P\le2$ is the polynomial power above. Thus
\[
 \sum_{\ell\in\mathbb Z^2}
 r_\ell^P\frac{W(Cr_\ell/L)}{W(cr_\ell)}\le C.
\]
Normalize the fine modes by their displayed bounds and absorb
their remaining factors into the numerical coefficients. This
proves the two asserted coefficient sums. Real and imaginary parts
may be separated at a fixed additional cost. Near a wall perform
the expansion on the reflected extension and average the modes
over the finitely many nearby reflections, with their prescribed
scalar or oriented-edge signs. Choose $\chi=1$ also on the required
image regions. The identities and normalized bounds persist.
The periodic boxes under consideration unwrap with the stated
clearance, so the same construction applies there.
\end{proof}

We next control the Gaussian source created by a collection of atoms.
The required estimate must cover every future gradient tail and every
point in an atom's neighborhood, uniformly in how many scales remain.
The two suprema in the following lemma express exactly these choices.

\begin{lemma}[Non-tangential square estimate]\label{xyh:maximal}
For a stacked coordinate vector $g=(g_r)_r$, put
$T_u g=\sum_{r\ge u}\nabla\Gamma_r^{1/2}g_r$, with binary scale
indices including the first and terminal pieces. For each fixed
aperture $A$ and each $q\ge0$,
\[
 \left\|\sup_u\sup_{|y-x|\le Au}
 u^q|\nabla^qT_ug(y)|\right\|_2
 \le Ce^{C(q+1)^2}
       \left(\sum_r\|g_r\|_2^2\right)^{1/2}.
\]
The same conclusion holds for the primal rotated gradient and for
their finite direct sums. By increasing $K_0$ once, the sum of these
maxima with weights $e^{-K_0q^2}$ has the same square bound.
All constants are uniform in depth and shell grouping, and the
conclusions include the reflected, periodic, and terminal pieces.
\end{lemma}
\begin{proof}
Write $\mathcal D_r=\nabla\Gamma_r^{1/2}$. The multiplier bounds in
Equations~\eqref{xyh:binary-multiplier} and
\eqref{xyh:high-derivatives} imply
\[
 |d_r(\xi)|\le C\min\{r|\xi|,
 e^{-c(\log^+(r|\xi|))^2}\}.
\]
Here $|\xi|$ may equivalently be replaced by the square root of the
lattice Laplacian symbol. On nonconstant terminal modes the same
bound holds because the terminal scale is comparable to the box
side. The white resolution, Equation~\eqref{xyh:white-resolution},
or Fourier Cauchy--Schwarz gives, for $F=\sum_r\mathcal D_rg_r$,
\[
 \|F\|_2^2\le C\sum_r\|g_r\|_2^2.
\]
Set $R_u=T_ug-Q_uF$. The term from scale $r$ has multiplier
$(1-Q_u)d_r$ for $r\ge u$, and $-Q_ud_r$ for $r<u$.
For every $q\ge0$,
\[
 \|(u\nabla)^qR_{u,r}\|_{2\to2}
 \le Ce^{C(q+1)^2}\min(r/u,u/r).
\]
For $r<u$, put $t=u|\xi|$. The multiplier is bounded by
$C(r/u)t^{q+1}e^{-c(\log^+t)^2}$. For $r\ge u$, put
$t=r|\xi|$ and $\rho=u/r$. Use
$|1-Q_u|\le C\min\{1,(\rho t)^2\}$ and
\[
 (\rho t)^q\min\{1,(\rho t)^2\}
 \le\rho^{q+1}t^{q+1}.
\]
The supremum of a fixed power of $t$ times the logarithmic Gaussian
decay costs at most $Ce^{C(q+1)^2}$. Young's Inequality on binary
scale indices consequently gives
\[
 \sum_u\|(u\nabla)^qR_u\|_2^2
 \le Ce^{C(q+1)^2}\sum_r\|g_r\|_2^2.
\]
For $q\ge1$ the same estimate holds directly for $T_ug$, because
only $r\ge u$ occurs and the derivative gives the summable gain
$(u/r)^q$.

The scaled lattice Sobolev estimate in two dimensions is
\[
 \sup_{|y-x|\le Au}|H(y)|^2
 \le C\sum_{k=0}^2u^{2k-2}
             \sum_{|z-x|\le Bu}|\nabla^kH(z)|^2.
\]
To obtain it, apply a cutoff on the larger box and use Fourier
Cauchy--Schwarz with weight $(1+u^2|\xi|^2)^2$; the inverse-weight
sum is $O(u^{-2})$. Summing the displayed estimate over $x$ and
then over $u$ proves the non-tangential bounds for $R_u$ and the
positive-order tails. At order zero the remaining term satisfies
\[
 \sup_{|y-x|\le Au}|Q_uF(y)|\le C M(|F|)(x),
\]
where $M$ is the lattice averaging maximal operator: $Q_u$ has
range $O(u)$ and kernel size $Cu^{-2}$. The square bound for $M$
follows from the cube covering inequality and layer-cake integration.
These arguments prove the asserted bounds. Their constants grow
at most as $Ce^{C(q+1)^2}$; summing with $e^{-K_0q^2}$ is legitimate
for sufficiently large fixed $K_0$.

Reflection reduces the finite wall problem to periodic extensions
with the required signs and a fixed volume multiplicity. The same
Fourier and local estimates hold there, counting norms per period.
The last remainder has the stated box-scale estimates. Harmonic
constants, where present, have zero approximation error since
$Q_u1=1$, zero positive differences, and a square bound from the
white resolution. This also covers the finite direct sums in the
statement. A grouped shell is obtained by restricting the binary
tail indices and by taking its corresponding coordinate direct
sum, which does not enlarge these estimates.
\end{proof}

Writing a sum of atom evaluations as a linear functional of the stacked
Gaussian coordinates produces a coefficient vector, the combined
source $I$ below.  The square estimate bounds its norm and the maximal
gradient tails produced by its scalar-coordinate part.  Here the assumptions
on birth-box overlap and scale weights are explicit; their validity for
each generated live root is checked in Lemma~\ref{xyh:history} and the
attenuation description that follows it.

\begin{lemma}[Atom source bounds]\label{xyh:atoms}
Consider $n_a$ normalized scalar or edge atoms, with a fixed bound
on the overlap multiplicity of their comparable birth boxes at
each birth size. Couple each atom only to field pieces at or after
its birth. Its scale coefficients may depend on the atom and have
uniformly bounded magnitude and total variation. The resulting
combined source on the stacked unmixed coordinates satisfies
\[
 \|I\|_2^2\le Cn_a.
\]
This holds on both coordinate components and their direct sum.
For its scalar-coordinate component, the output tails
$v^{\ge u}=\sum_{r\ge u}D_rI_r$ satisfy
\[
 \left\|\sup_u\sup_{|y-x|\le Au}|v^{\ge u}(y)|\right\|_2^2
 \le Cn_a.
\]
\end{lemma}
\begin{proof}
Let $\mathcal M g$ be the derivative-weighted maximal function in
Lemma~\ref{xyh:maximal}. A scalar atom born at size $h$ may subtract
the future scalar field's anchor value. A lattice path bounds the
zeroth order by $Ch$ times its gradient maximum; the higher test
orders are controlled by the corresponding terms in $\mathcal M$.
An edge atom obeys the same bound because its test normalization
contains the extra factor $h$. Thus either atom acts on future
fields with absolute value at most
\[
 Ch\inf_{x\in B}\mathcal M g(x),
\]
where $B$ is its comparable birth box and the fixed aperture is
large enough to contain the required collar. For each atom Abel
summation writes a weighted future sum exactly as
\[
 \sum_{r\ge h}a_rF_r
 =a_hT_h+\sum_{r>h}(a_r-a_{r^-})T_r.
\]
The sum of coefficient magnitudes is uniformly bounded by the
assumed magnitude and total variation, so the preceding atom bound
persists. All tails here start at a size at least the atom birth.

For the family of birth boxes put
\[
 \rho(x)=\sum_\alpha h_\alpha^{-1}1_{B_\alpha}(x).
\]
Bounded overlap and geometric summation over dyadic sizes give
\[
 \|\rho\|_\infty\le C,\qquad
 \rho(x)^2\le C\sum_\alpha h_\alpha^{-2}1_{B_\alpha}(x),
 \qquad \|\rho\|_2^2\le Cn_a.
\]
For the middle inequality compare both sides to the inverse square
of the smallest occupied size at $x$. Since $|B_\alpha|\asymp
h_\alpha^2$, duality and Lemma~\ref{xyh:maximal} give
\[
 \begin{split}
 |\langle I,g\rangle|
 &\le C\sum_\alpha h_\alpha
                 \inf_{B_\alpha}\mathcal M g
 \le C\sum_x\rho(x)\mathcal M g(x)\\
 &\le C\sqrt{n_a}
       \left(\sum_r\|g_r\|_2^2\right)^{1/2}.
 \end{split}
\]
This proves the source bound on all stated components. Apply
Lemma~\ref{xyh:maximal} once more, now to the scalar-coordinate
source itself, to obtain the output square bound. The proof uses
the finite extensions and terminal estimates already included in
that lemma, so it applies to all the stated geometries.
\end{proof}

The source square bound pays the quadratic costs of Gaussian
completion.  The branch bound in Equation~\eqref{xyh:primitive-bound}
also contains exponential growth in the imaginary argument.
We therefore need all moments of the resulting output fields.
The next lemma uses their spatial decay and the John--Nirenberg
estimate to obtain the moment growth used later in
Equation~\eqref{xyh:imaginary-payment}.

\begin{lemma}[All moments of atom outputs]\label{xyh:moment-bound}
Under the hypotheses of Lemma~\ref{xyh:atoms}, let
$v(x)=\sum_r b_r(x)D_rI_r(x)$, where at each output point the
coefficients $b_r(x)$ have uniformly bounded magnitude and total
variation in scale. They need not have spatial regularity. Then
\begin{equation}\label{xyh:atom-moments}
 \sum_x|v(x)|^p\le n_a(Cp)^p,\qquad p\ge2.
\end{equation}
The sum is over the physical lattice, or over one period in a
periodic extension, with the fixed reflection multiplicity
absorbed into $C$. The constants are uniform in depth and shell
grouping.
\end{lemma}
\begin{proof}
First omit the output coefficients and write
$f=\sum_\alpha f_\alpha$ for the full output. Keep the allowed
atom-dependent input coefficients. Testing a binary covariance
derivative kernel at $r\ge\max(h_\alpha,u)$ against its atom gives
$Ch_\alpha r^{-2-j}$ for $j=0,1$. For scalar atoms subtract a kernel
constant before testing; for edge atoms the second covariance
derivative and edge normalization give exactly the same bound.
The high derivative bounds in Equation~\eqref{xyh:high-derivatives}
are absorbed by the test weights, using $K_0$ as fixed above.
The kernel vanishes unless its range reaches the atom box.
Summing over binary pieces, and using the uniform bound on the
input coefficients, yields
\begin{equation}\label{xyh:single-atom-tail}
 |\nabla^jf_\alpha^{\ge u}(x)|
 \le\frac{Ch_\alpha}
 {(h_\alpha+u+\operatorname{dist}(x,B_\alpha))^{2+j}},
 \qquad j=0,1.
\end{equation}
The terminal piece satisfies the same estimate at box scale.

Let $Q$ have side $u\ge1$. Call an atom near if $h_\alpha\le C_0u$
and its birth box is within $C_0u$ of $Q$; all others are far.
Choose $C_0$ greater than the range and cone constants. The number
of near atoms at size $h$ is at most $C(u/h)^2$, and their total
number is at most $Cu^2$. Lemma~\ref{xyh:atoms}, applied just to
this subfamily and then to the full output operator, gives
\[
 \frac1{|Q|}\sum_Q|f_{\rm near}|\le C.
\]
Equation~\eqref{xyh:single-atom-tail} bounds their tail from $u$
pointwise on $Q$ by
\[
 \sum_{\alpha\ {\rm near}}
       \frac{Ch_\alpha}{(h_\alpha+u)^2}
 \le C\sum_{\substack{h\le C_0u\\h\ {\rm dyadic}}}h^{-1}
 \le C.
\]
On $Q$ the full output of each far atom equals its tail from $u$:
its birth size is larger, or finite range excludes all smaller
pieces. The variation of all far outputs across $Q$ is at most
\[
 C\sum_\alpha
 \frac{u h_\alpha}
 {(u+h_\alpha+\operatorname{dist}(Q,B_\alpha))^3}
 \le C\sum_{h\ {\rm dyadic}}\frac{u}{h(u+h)}\le C.
\]
For the middle bound, bounded overlap at size $h$ compares the
spatial sum with $h^{-2}$ times the integral of the kernel;
that integral is $O((u+h)^{-1})$. Small sizes contribute
$C\sum_{h\le u}h^{-1}$ and large ones
$C\sum_{h>u}u h^{-2}$, both bounded. It follows that
\[
 \frac1{|Q|}\sum_Q|f-f_Q|\le C,
 \qquad |f^{\ge u}(y)-f_Q|\le C\quad(y\in Q).
\]
In particular $f$ has uniformly bounded BMO norm and, with a fixed
enlargement of $Q$ for the cone,
\[
 T_*I:=\sup_u\sup_{|y-x|\le Au}|f^{\ge u}(y)|
 \le C M(|f|)+C.
\]
The same kernel estimates bound unit-scale neighbor differences.
Thus the piecewise-constant extension has bounded ordinary cube
BMO norm, including cubes smaller than a lattice cell.

Here is the needed exponential BMO estimate, in the form of the
John--Nirenberg Inequality~\citep{JohnNirenberg1961}. In a cube $Q$, stop
on maximal dyadic subcubes where the average of $|f-f_Q|$ exceeds
twice the uniform BMO bound. Their total volume is at most
$|Q|/2$. Their averages differ from $f_Q$ by a fixed multiple of
that bound, because their parents did not stop. Repeat inside each
selected cube with its own average. After $k$ generations the
selected union has volume at most $2^{-k}|Q|$; outside it,
$|f-f_Q|\le C(k+1)$. Consequently
\[
 |\{x\in Q:|f(x)-f_Q|>t\}|\le C|Q|e^{-ct}.
\]
This argument applies componentwise to real and imaginary parts.

Select maximal averaging cubes of $|f|$ above a sufficiently large
fixed level. By Lemma~\ref{xyh:atoms} and the square maximal
inequality their total volume is $O(n_a)$. The averages on these
cubes are uniformly bounded, since their parents are below the
chosen level. Outside their union $|f|$ is below that level.
Applying the preceding exponential estimate on the selected cubes
therefore gives
\[
 |\{x:|f(x)|>t\}|\le Cn_a e^{-ct},\qquad t\ge C.
\]
Below this fixed level use $\|f\|_2^2\le Cn_a$ instead.
Layer-cake integration gives
\[
 \sum_x|f(x)|^p\le n_a(Cp)^p,\qquad p\ge2.
\]
For completeness the $L^p$ maximal bounds used here need no
endpoint hypothesis: the cube covering inequality applied to
$|f|1_{\{|f|>t/2\}}$ bounds $|\{M|f|>t\}|$ by
$Ct^{-1}\sum_{|f|>t/2}|f|$; integration gives a constant uniformly
bounded in operator norm for $p\ge2$.

The additive constant in the comparison of $T_*I$ with $M(|f|)$
is not summed over the whole plane. On $\{T_*I\le C_1\}$ use
$(T_*I)^p\le C_1^{p-2}(T_*I)^2$ and the separate square bound in
Lemma~\ref{xyh:atoms}. On its complement choose $C_1$ large enough
that $T_*I\le C M(|f|)$ and use the $L^p$ maximal inequality.
This proves $\sum_x(T_*I)^p\le n_a(Cp)^p$.
Finally Abel summation at each output point bounds the requested
weighted output by $C T_*I$, proving
Equation~\eqref{xyh:atom-moments}.

In finite reflected or periodic geometries, perform the argument
on the prescribed extensions and count volume per period.
Bounded overlap implies that the number of atoms per period is
at most a constant times its volume; the square estimate therefore
bounds the whole-period average by a fixed constant. Choose the
stopping level above it. All stopping cubes then have the stated
parent bound, and the same proof gives the required exponential
tail with prefactor $Cn_a$. The final box-scale remainder is
included in Equation~\eqref{xyh:single-atom-tail}; harmonic
constants, if present, have no variation and obey their square
bound. Thus neither the finite endpoint nor reflection changes
any constant in the conclusion.
\end{proof}

\subsection{The precise class of generated histories}

We now specify the expansion records described above.  The point of
the rules is to ensure that repeated local operations preserve the
support separation and bounded overlap required by the atom estimates.
They also determine which microscopic charges remain coupled to future
Gaussian fields, information needed for the separate charge budget.

Use the connected-activity algebra of Proposition~\ref{rg:partition-map}
with field $z=(w,\psi)$. Shift sectors refer to uniform translations of
$\psi$ with $w$ fixed, and evenness to simultaneous negation. For an
$m$-block $B$ write
\begin{equation}\label{xyh:tensor}
 \mathcal E_B(t)=\frac12\sum_e\mathrm{wt}_B(e)
 \{t_{11}w_e^2+2i t_{12}w_e(\partial\psi)_e
                         +t_{22}(\partial\psi)_e^2\},
\end{equation}
where the endpoint weights assign half of each incident edge to its block.
A tensor primitive is this polynomial with an already determined numerical
coefficient. Its Gaussian expectation is the same polynomial plus a
constant; assign the whole polynomial to the coarse singleton and only
the contracted numerical constant to the scalar normalizer. The
background translate of a gradient polynomial is never part of that
normalizer.

On a small bulk support the neutral even remainder is localized by its
constant and quadratic Taylor coefficients. Evaluate the Hessian on
scalar affine directions with slopes $m^{-1}$ and on constant oriented
edge directions of size $m^{-1}$. Contract with
$(m w(v),m\nabla\psi(v))$, average the choices of signed orientation,
and then average the anchors in the fine blocks. This is an exact
prescription. Translation and square symmetries make the summed
coefficients a tensor of the form Equation~\eqref{xyh:tensor}. Complex
conjugation on real fields changes $\psi$ to $-\psi$, so its three
parameters are real. At an even wall subtract the neutral constant;
at a defect subtract the neutral constant without an evenness assumption.
Good-class copies use the same prescriptions on neighborhoods where the
kernels agree. Local source translations act on both fields according to
Equation~\eqref{xyh:offsets} and on a local lift for the scalar.

At the microscopic start expand by selecting nonzero vertex charges,
edge factors $F_b$, and quadratic factors $e^{U_e}-1$. Use the exact identity
\begin{equation}\label{xyh:quadratic-remainder}
 e^{U_e}-1=U_e+\int_0^1(1-\tau)U_e^2e^{\tau U_e}\,d\tau.
\end{equation}
Both endpoints belong to an edge factor's support. After the first
independent piece and temporary grouping, relocate the complete linear
$U_e$ expectation with half endpoint weights, keeping its original
exclusion. Its contraction is a designated scalar independent of the
source offsets; the remaining shifted polynomial is retained as a
polynomial. The initial bulk tensor is
$t_0=(s_{11},s_{12},s_{22})$.

\begin{definition}[A selected history]\label{xyh:history-definition}
First truncate the microscopic charge indices and use strict compensation
as in Proposition~\ref{xyh:compensation}. Expand every finite sum,
Taylor integral, Fourier direction choice, vertex/orientation average,
and Cauchy contour, and fix one term. Cut its dependency graph at every
scalar coefficient evaluation. The resulting record is a rooted forest
of separately completed Gaussian roots and at most one live Gaussian
root. The forest is understood by occurrences: if the same evaluated
coefficient is used twice, its completed calculation has two separate
copies when displayed. A live root is a maximal Gaussian calculation not
yet completed to a scalar. Its shell variables and any future field pieces
are represented on one coupled ambient stack. Its events stay in that root
until completion, even if a current attenuation becomes zero. A completed
root has already been evaluated to a number and is never put
on that stack again. The leaves of a live root are selected microscopic
factors, new polynomial primitives, and opaque occurrences of previously
completed numerical coefficients. Its internal vertices are the following
operations:
\begin{enumerate}[label=\textup{(\roman*)}]
\item A convolution of a compatible family of fine components with one
shell, followed by temporary reblocking. A term with two or more inputs
has no Taylor interpolation at that step.
\item A small linear, single-input localization remainder. Its integral
Taylor formula evaluates its one child at $\rho z$, $0\le\rho\le1$,
with the same real $\rho$ in both fields. Each Taylor direction is first
expanded using Lemma~\ref{xyh:fourier}; its atom value is kept as a
separate factor and its normalized child direction is fixed before
Cauchy differentiation.
\item A scalar coefficient calculation, which completes its child root.
An assigned constant or polynomial then uses an opaque occurrence of
that evaluated number as a coefficient. The polynomial's new linear
factors start in the new live root and have no microscopic leaves or
atoms of the completed root attached to them.
\item A product of distinct temporary objects obeying their original
exclusions, regrouped by final connectivity, or a product with scalar
normalizer decorations. Splitting off a free singleton tensor is a
choice of a summand, not a duplication of its previous history.
\end{enumerate}
A leaf is called live if its dependence on future Gaussian fields remains
in the selected term. The numbers $N_h,n_F,n_{\rm ch}$ and the list of
original charges of a root refer only to that root. Here $N_h$ counts its
microscopic leaves and actual polynomial, Taylor, and probe events;
fixed-degree contour representations add a bounded number of events per
operation. A product that joins live components joins them into one live
root, using the common-stack construction below. A product with an opaque
coefficient occurrence adds no live Gaussian coordinate or atom.

The coefficient dependencies are acyclic because a coefficient is used
only after its evaluation. Completed numbers may be combined numerically
before a new primitive is formed; this uses the resulting number and does
not revive any of their Gaussian roots. If a completed coefficient $a$ is
used $m$ times, there are $m$ opaque occurrences. Expanding those occurrences for
bookkeeping gives $m$ separately completed copies, not one coupled root.
In particular an inverse-normalizer term $a^m$ has the product of $m$
completed-root bounds. Every overlap, source-square, and charge-payment
assertion below is per live root; after final evaluation it applies to
each completed root separately. The weight for the full forest is their
recursive product, as specified in Proposition~\ref{xyh:calculus}.
\end{definition}

Cauchy's formula here is an identity for a derivative of one shifted copy
of its child, not a product of copies. For example, after writing
$z=\sum_\ell c_\ell A_\ell(z)u_\ell$ on the needed region,
a Taylor argument contributes $c_\ell A_\ell(z)$ and
$D_{u_\ell}H(\rho z)$ in a selected term. Only the fixed $u_\ell$ enters
the Cauchy shift of $H$. An atom factor can itself be written as the
contour derivative of $e^{tA_\ell(z)}$ on a fixed circle. A later
Taylor operation applies the same rule to the resulting product.
Consequently it adds deterministic probes and real attenuations;
it never multiplies an original charge by a complex interpolation
parameter. An original phase charge on a Gaussian field remains real.

Figure~\ref{xyh:history-figure} separates three operations and records
what happens to their live microscopic ancestry.

\begin{figure}[tb]
\centering
\begin{tikzpicture}[
  x=1cm,y=1cm,
  font=\footnotesize,
  history box/.style={draw=black!55,rounded corners=2pt,fill=white,
    align=center,inner sep=4pt},
  history live/.style={history box,fill=blue!4},
  history edge/.style={-{Stealth[length=4pt]},draw=black!70,line width=.55pt},
  history note/.style={align=center,font=\scriptsize},
  history title/.style={font=\footnotesize\bfseries,align=center}
]
  \foreach \x in {0,5.36,10.72}
    \draw[draw=black!18,fill=black!2,rounded corners=3pt]
      (\x,0) rectangle ++(5.08,5.55);

  \node[history title] at (2.54,5.27) {(a) Single-input Taylor};
  \node[history live,text width=4.28cm] (htout) at (2.54,4.45)
    {$c_\ell A_\ell(z)\,D_{u_\ell}H(\rho z)$\\[2pt]
     \scriptsize new atom $A_\ell(z)$\\[-1pt]
     \scriptsize or its contour probe};
  \node[history live,text width=4.00cm] (htin) at (2.54,1.16)
    {one live child $H(z)$\\[1pt]
     \scriptsize microscopic leaves on $X$};
  \draw[history edge] (htin.north) -- (htout.south)
    node[midway,fill=black!2,align=center,inner sep=3pt]
      {$\rho z=(\rho w,\rho\psi)$\\[-1pt]
       \scriptsize same real $\rho\in[0,1]$ in both fields};
  \node[history note,text width=4.70cm] at (2.54,.30)
    {Cauchy: $H(\rho z+t u_\ell)$ once.\\
     Support: $X\subseteq X'$.};

  \node[history title] at (7.90,5.27) {(b) Two-input merger};
  \node[history live,text width=4.00cm] (hmout) at (7.90,4.45)
    {merged live history\\[-1pt]
     \scriptsize support $X'$};
  \node[history box,text width=4.00cm,font=\scriptsize] (hm) at (7.90,2.81)
    {one shared current shell\\[-1pt]
     convolution\\[-1pt]
     temporary reblocking\\[-1pt]
     \bfseries no Taylor at this vertex};
  \node[history live,text width=1.61cm] (hminone) at (6.68,1.12)
    {$H_1$\\[-1pt]\scriptsize on $X_1$};
  \node[history live,text width=1.61cm] (hmintwo) at (9.12,1.12)
    {$H_2$\\[-1pt]\scriptsize on $X_2$};
  \draw[history edge] (hminone.north) -- ([xshift=-23pt]hm.south);
  \draw[history edge] (hmintwo.north) -- ([xshift=23pt]hm.south);
  \draw[history edge] (hm.north) -- (hmout.south);
  \node[history note] at (7.90,.57) {compatible inputs};
  \node[history note] at (7.90,.20) {$X_1\cup X_2\subseteq X'$};

  \node[history title] at (13.26,5.27) {(c) Coefficient detachment};
  \node[history live,text width=4.28cm] (hdnew) at (13.26,4.45)
    {\textbf{new} polynomial primitive\\[1pt]
     $aL_1(z)L_2(z)$\\[1pt]
     \scriptsize fresh linear factors $L_1,L_2$};
  \node[history box,text width=3.72cm] (hdeval) at (13.26,2.74)
    {scalar evaluation: $a$\\[1pt]
     \scriptsize\bfseries old child history ends};
  \node[history box,text width=4.00cm] (hdold) at (13.26,1.16)
    {old child $H_{\rm old}$\\[1pt]
     \scriptsize microscopic leaves\\[-1pt]
     \scriptsize and older atoms};
  \draw[history edge] (hdold.north) -- (hdeval.south);
  \draw[history edge,dashed] (hdeval.north) -- (hdnew.south)
    node[midway,right,fill=black!2,inner sep=2pt,font=\scriptsize]
      {number only};
  \node[history note,text width=4.70cm] at (13.26,.25)
    {No old leaves attached to $L_1,L_2$.};
\end{tikzpicture}
\caption{Selected history operations from
Definition~\ref{xyh:history-definition} and Lemma~\ref{xyh:history}.
These are separate schematic rules; panel~(a) shows one selected Taylor
direction factor. Solid arrows carry live field dependence up to its
termination; the dashed arrow carries only an evaluated number.
Live microscopic supports only enlarge or unite, Cauchy differentiation
uses one shifted child, and detachment ends the old microscopic and atom
ancestry before new polynomial factors are formed.}
\label{xyh:history-figure}
\end{figure}

\begin{lemma}[Supports, collars, and overlap]\label{xyh:history}
Each live Gaussian root of Definition~\ref{xyh:history-definition} has
the following uniform properties. They also apply when that root is
completed. No assertion combines atom boxes or source squares of distinct
completed roots in the same forest.
\begin{enumerate}[label=\textup{(\alph*)}]
\item A support carrying microscopic leaves or atoms of the live root never shrinks: it is
replaced only by coarse closure or a union. No selected microscopic
factor of the same type occurs twice in one live root.
\item All fields in the live root can be represented on one ambient
stack of shell coordinates. Individual atom representations require a
fixed collar, whose width in current block units is bounded independently
of depth and $L$.
\item Atom boxes have bounded overlap at each birth scale, uniformly in
$L$ and history depth. The same holds for the fine-scale probe boxes of
one step.
\item If a Taylor operation after the shell ending at $S$ changes the
attenuation on one group of live microscopic leaves, then groups having
different relative changes have zero cross entries in the restricted
white tail $E_{\ge S}$.
\end{enumerate}
All assertions apply to an early-stopped root, reflected grids, and
the auxiliary torus.
\end{lemma}
\begin{proof}
Initially each selected vertex factor appears once at its vertex, and each
selected edge factor of a specified type appears once and contains its two
endpoints. Convolution takes closures, a residual keeps
the full temporary support, and products take unions. The only shrinking
operation is coefficient detachment in Definition~\ref{xyh:history-definition}(iii),
where the old Gaussian root, including its atoms, has ended. A new
polynomial begins with its own support and atoms. Scalar decorations
introduce no field, and choosing a summand or one contour derivative does
not copy a microscopic factor. This proves (a) inductively within each
root. A tilt factor and an
$F_b$ factor may use the same physical edge; the no-repetition assertion
is separate for each type, as required below.

For (d), fix a Taylor change at scale $S$. Its child is a single input, so
all live leaves in that child receive the same relative factor. For any
other leaf with a different factor, follow the branches to their first
product ancestor. If that product is formed at this regrouping, they
were distinct temporary objects and their original supports have
distance greater than $r_0S$. If it occurs later, the two unshrunk
ancestors are distinct compatible fine components, or distinct temporary
objects at that later step. Their distance is therefore greater than
the corresponding fixed radius times a scale at least $S$. There is no
case of differently changed live branches inside the same temporary
Taylor input: that input was linear. By (a), later closure cannot have
hidden an earlier nearby leaf. Choose $r_0$ larger than the range constant
in Equation~\eqref{xyh:white-resolution}. Its white tail then has the
asserted zero cross blocks. Unchanged leaves count as a group with
relative multiplier one.

Here is the collar induction needed for (b). If a newly emitted atom uses
$c_0$ birth-block widths and all inherited dependence uses $c_*$ old-block
widths, the new dependence is contained in
$c_0+c_*/L+c_{\rm st}$ current-block widths, with a fixed stencil allowance
$c_{\rm st}$. Fix $c_*=2(c_0+c_{\rm st})+1$; the required induction
inequality then holds for every $L\ge2$. Let $c_{\rm ker}$ bound the
relevant kernel range in current block units. The small-closure conclusion
of Lemma~\ref{rg:geometry} makes a support relocation at most one current
block, so fix $c_{\rm move}=1$ for that conclusion. Choose the temporary radius
$r_0>2c_*+c_{\rm ker}+4$ and the compatibility radius
$R>r_0+2c_{\rm move}+4$, increasing them if necessary for earlier
kernel ranges. The small-support threshold in Lemma~\ref{rg:geometry}
uses these fixed radii, and $L$ is then taken large enough for its
small-closure conclusion. The resulting fine small-support diameter is
fixed before this last choice of $L$, so the Fourier cutoff uses a fixed
number of coarse blocks. These separations
are the locality hypotheses used in the factorization proof of
Proposition~\ref{rg:partition-map}, with the present block collars in
place of its smaller microscopic collar. Before
compatible inputs meet, all their already integrated shells then have
zero cross covariance, including their atom dependence. The current
shell at a multiple-input convolution is shared from the outset.
Before products of temporary objects, use their stronger temporary
separation. Gaussian characteristic functions factor when these cross
blocks vanish; Fourier uniqueness proves factorization for the finite
complex Gaussian measures of the strict regularization. The limiting
factorization follows after the endpoint bounds below. This proves the
common-stack induction. When a finite-grid good piece is copied from the
plane, only equality of the restricted covariances is needed; the chosen
square roots need not be identical.

For (c), a selected Taylor term emits only its fixed number of direction
atoms at the upper birth scale, and only a single-input temporary object
emits them. Distinct batches in one final history are separated at that
birth scale by the same first-product-ancestor argument. In particular
an upper block may contain $O(L^2)$ fine inputs, but its multiple-input
convolution emits no new Taylor atoms: the old atoms retain their old
birth scales. A tensor is a normalized vertex/orientation sum; a selected
summand contains only its two linear factors, not a pair at every vertex.
The same observation applies to a localized quadratic assignment.
Products of these primitives have fine compatibility or temporary
separation. At a fixed scale there are only finitely many such types.
If a coefficient calculation ends, its older atoms end with it and are
not stacked on its newly assigned polynomial. Repeated opaque occurrences
of that coefficient add no atoms to the new root; if expanded, each copy
is checked on its own completed stack. At the microscopic scale
the list of primitive types and polynomial degrees is fixed. This proves
the bounded overlap without an $L^2$ loss. Probes have the same ancestry
and a fixed number per node. Reflections, periodic distance, and the
terminal harmonic piece do not alter any of these support arguments.
\end{proof}

For one shell introduce coordinates $x=(x_T,x_L),y$ with inverse
covariance
\[
 \begin{pmatrix}\mu I&i d_cJ\\i d_cJ^t&sI\end{pmatrix},
 \qquad Jy=(0,y),
\]
and fields $w_r=B_rx_T+D_rx_L$, $\psi_r=\Gamma_r^{1/2}y$,
$p_r=D_ry$. Use the same letter $J$ for the orthogonal embedding on the
complete stack. Each microscopic leaf and each atom has attenuation
one at birth; passing a Taylor node multiplies its future coupling in
both fields by the same number in $[0,1]$. It may eventually terminate.
Thus these scale weights have uniformly bounded variation.

\begin{lemma}[Contraction of the attenuated stack]\label{xyh:contraction}
Fix one live Gaussian root. On its distinct microscopic tilt edges, let
$R_wx,R_py$ be the attenuated
edge and gradient totals. On the distinct $F_b$ edges let $R_Fx$ be its
edge total. Then
\begin{equation}\label{xyh:contracting-maps}
 \|R_w\|\le1,\qquad\|R_p\|\le1,\qquad
 \|R_F\|\le1,\qquad R_wJ=R_p.
\end{equation}
\end{lemma}
\begin{proof}
Restrict the white kernels to the distinct edges under consideration and
write $A_r$ for their diagonal attenuations. The row covariance is
$\sum_rA_rE_rA_r$. If the history stops early, add the missing positive
tail with the last attenuation. Remove changes in reverse chronological
order. After all later changes have been removed, a change at $S$ leaves
a tail $P Q E_{\ge S}Q P$, where $P$ records all earlier factors and
$Q$ the new relative factors. By Lemma~\ref{xyh:history}(d), the restricted
tail is block diagonal between groups on which $Q$ differs. On each
block $Q=qI$ with $0\le q\le1$, whence
$QE_{\ge S}Q\preceq E_{\ge S}$. Congruence by $P$ preserves this order,
even when $P$ varies within a block. Removing every change leaves
$\sum_rE_r=I$. This proves the $R_w$ and $R_F$ bounds. Replacing
$D_rD_r^*$ by $E_r$ proves the $R_p$ bound. Finally the same attenuation
in both fields gives $R_wJ=R_p$ row by row. Deterministic probe and twist
backgrounds do not change these linear maps.
\end{proof}

\subsection{Finite Gaussian completion and endpoint bounds}

\begin{lemma}[Mixed completion and its Cameron bound]\label{xyh:mixed}
Fix one live Gaussian root. After incorporating its microscopic quadratic
tilts, let $\Theta$ be their diagonal matrix of strengths in $[0,1]$.
The stacked inverse covariance has blocks
\begin{equation}\label{xyh:completed-precision}
 \begin{gathered}
 \begin{pmatrix}M&iH\\iH^t&S_0\end{pmatrix},\qquad\text{where}\\
 \begin{aligned}
 M&=\mu I-s_{11}R_w^t\Theta R_w,\\
 H&=d_cJ-s_{12}R_w^t\Theta R_p,\\
 S_0&=s(I-R_p^t\Theta R_p)\succeq0.
 \end{aligned}
 \end{gathered}
\end{equation}
Its Gaussian reweighting determinant has modulus at most $e^{CN_h}$.
Define
\begin{equation}\label{xyh:completion-matrices}
 D=(S_0+H^tM^{-1}H)^{-1},\quad C'=M^{-1}HD,
 \quad V=M^{-1}-C'D^{-1}(C')^t\succeq0.
\end{equation}
Then $M,D$ are uniformly close to the identity. A source $e^{i(Y,y)}$
contributes $e^{-(Y,DY)/2}$, and the remaining $x$ integral is a real
Gaussian with covariance $V$ evaluated at mean $C'Y$.
If $I=\Im Y$, the imaginary evaluation shift $iC'I$ may be replaced by
$iJI$ at cost at most $e^{C\|I\|^2}$, leaving the remaining real
Gaussian estimate unchanged.
\end{lemma}
\begin{proof}
Equation~\eqref{xyh:completed-precision} follows by substituting the
linear maps of Lemma~\ref{xyh:contraction} into the quadratic tilts.
Those contraction bounds imply $M=I+O(|s_\bullet|)$ and
$H=J+O(|s_\bullet|)$, while $0\preceq S_0\preceq sI$.
The Schur complement is therefore positive and close to $I$.
The perturbation of the untilted inverse covariance has rank at most a
fixed multiple of the number of tilt leaves. The untilted covariance has
bounded operator norm, including as $s\downarrow0$; its product with the
precision perturbation is small in operator norm. Thus the determinant
ratio is a product of at most $CN_h$ numbers of modulus in a fixed
neighborhood of one. This proves its bound independently of the number
of unused Gaussian coordinates.

For $S_0>0$, integration in $y$ or direct block inversion gives covariance
\[
 \begin{pmatrix}V&-iC'\\-i(C')^t&D\end{pmatrix}.
\]
Its characteristic function with a $y$ source proves the asserted mean
and exponential factor. Positivity of $V$ is the Schur-complement
identity
$V^{-1}=M+HS_0^{-1}H^t$; limits prove the positive semidefinite case.

Let $A=I-R_p^t\Theta R_p$, $\Delta=C'-J$. Using $R_wJ=R_p$ gives
\begin{align*}
 H^t\Delta
 &= (I-S_0D)-(d_cI-s_{12}R_p^t\Theta R_p)\\
 &= -A(sD+s_{12}I).
\end{align*}
For nonsingular $A$ and $s>0$ the Cameron square satisfies
\[
 \Delta^tV^{-1}\Delta
 =\Delta^tM\Delta+
 (sD+s_{12}I)^t A(sD+s_{12}I)/s\preceq CI,
\]
since $s_{12}^2/s\le d_*^{-1}$. The same inequality holds on the
supported space by approximation when $A$ is singular. At $s=0$ one has
$s_{12}=0$ and $H^t\Delta=0$, which is exactly the required range
condition for the possibly singular $V$.

Translate the real Gaussian variable by $i\Delta I$. Its density changes
by a phase of modulus one and the factor
$\exp(\|\Delta I\|_{V^{-1}}^2/2)\le e^{C\|I\|^2}$.
The remaining evaluation shift is $iJI$. The entire-function translation
is justified first with positive covariance and strict compensation by
Equation~\eqref{xyh:primitive-bound}, then by the same Gaussian majorant
on the supported limiting space. Linear $x$-source factors can be left
in the integrand during this translation; their extra deterministic
factor is bounded by Cauchy--Schwarz using their source norm. No
unrestricted imaginary Gaussian is substituted into $F_b$.
\end{proof}

We next specify the deterministic decorations retained inside a live-root
integral. A polynomial contour has fixed radius and its atom factor has
the form $e^{t[A(Z)+d_A]}$, where $Z$ is its attenuated variable field.
The deterministic value $d_A$ is part of this exponential. Likewise,
deterministic offsets remain in the arguments of $F_b$ and of the
microscopic quadratic tilts. These decorations are not structural
numerical coefficients, even after the contour variables have been fixed.
Only the following rules introduce them: the fixed Cauchy probes in a
selected Taylor calculation; the intrinsic microscopic offsets in
Equation~\eqref{xyh:offsets}; and one good-support transport defined
below. They cannot contain an arbitrary numerical multiplier.

Every retained event of a live root carries at most one physical source
decoration. Each root also has a tag state. The unmarked construction
omits the physical-source record, so its roots start untagged. A present
intrinsic physical-source record tags its root at birth even when all its
numerical offsets vanish. A good-support transport is defined only on an
untagged root and makes it tagged. The
empty root stays untagged. A representation is untagged when every
selected live root is untagged. A product root is tagged if any input root is
tagged, and Taylor operations never clear that state, even when a current
attenuation becomes zero. The transport leaves every opaque completed
coefficient occurrence untouched. Subsequent convolution and Taylor
operations carry existing physical decorations without applying that
transport again; new Taylor atoms are outside the differentiated child
and do not receive its same-step probe. Completion ends all the root's
decorations and its tag state. A new polynomial primitive with a completed
numerical coefficient starts with new atoms and its own untagged state.

For the early class there is a fixed constant $C_{\rm tag}$ such that the
total deterministic offset at each selected microscopic edge and $|d_A|$
at each atom are at most $C_{\rm tag}$. Physical scalar phase decorations
are real, so their exponentials have modulus one. The constant is uniform
for at most two signed sources, and for the early part of each fixed
separated source configuration specified below. This class also includes
the accumulated Cauchy probe decorations. Their uniform bound follows
directly from the birth normalization. If $A$ was born at $h$ and $u$ is
a later normalized probe at scale $m\ge h$, then
\[
 |A(u)|\le C h/m.
\]
For scalar atoms subtract the probe's anchor value and use constant
annihilation; for edge atoms use the extra factor $h$ in the test norm.
There are a fixed number of probes per passing scale, and
$\sum_{m\ge h,\ m/h\text{ a power of }L}h/m\le(1-L^{-1})^{-1}$.
At microscopic edges the corresponding sum is bounded by
$C\sum_{m\ge1}m^{-1}$ over the geometric scales. Thus these rules give
the asserted fixed bounds. The exception is a probe paired with an
original charge phase. At a Taylor node with current charge array $k$
this contributes at most
\begin{equation}\label{xyh:phase-probe-cost}
 \exp\{C h_m(k;X)\}.
\end{equation}
These factors are retained separately and paid in Lemma~\ref{xyh:budget}.

Here are the bounds for the physical decorations. On a good support write
$\Lambda=(\Lambda_w,\Lambda_\psi)$, where
$\Lambda_w=-i d_c\eta/\mu$ and the real local lift satisfies
$\partial\Lambda_\psi=\eta$ and agrees with $\vartheta$ modulo its period.
Let $A$ be an atom born at $h$, and let $\mathcal C_A$ be its fixed
expanded birth collar. Write $v^\flat$ for the physical representative of
$v$ under the prescribed wall reflections, with $v^\flat=v$ on the physical
grid. At every point $v\in\mathcal C_A$, the high-order source bound is
\[
 |\nabla^j\eta(v)|\le C e^{C_0(j+1)^2}b^{-1/2}
       \sum_i(1+\operatorname{dist}(v^\flat,x_i))^{-1-j}.
\]
This is a pointwise bound on the collar; it does not choose one source
as nearest throughout that collar. Put
$d_i(A)=\inf_{v\in\mathcal C_A}\operatorname{dist}(v^\flat,x_i)$.
When the collar has good clearance, the test normalization and the bound just displayed give
\[
 |A(\Lambda)|\le Cb^{-1/2}\sum_i\frac{h}{1+d_i(A)}.
\]
For a scalar atom, its zeroth order is a path of length $Ch$ times the
gradient bound; its order-$j$ contribution from source $i$ is bounded by
$C b^{-1/2}e^{-(K_0-C_0)j^2}(h/(1+d_i(A)))^j$. For an edge atom the same
calculation uses $h^{j+1}\nabla^j\eta$. The good clearance bounds
$h/(1+d_i(A))$ by a fixed constant, so the higher orders are absorbed
by the test weights. In particular, suppose the current support has
clearance $Hm$ and its inherited collars extend at most $c_{\rm col}m$
from its hull, with $H>c_{\rm col}$. Then
$d_i(A)\ge(H-c_{\rm col})m$. For one or two sources the bound for an old
atom of birth size $h\le m$ is therefore $Ch/m$.

For $k_{\rm src}$ sources at mutual distance at least
$\epsilon M_{\rm box}$, call an atom early when its own birth size obeys
\begin{equation}\label{xyh:early-source-scales}
 Lh\le c\epsilon M_{\rm box}/(1+k_{\rm src}).
\end{equation}
Choose $c$ so the birth collar has diameter less than
$\epsilon M_{\rm box}/4$. At each folded point of that collar at most one source
is within $\epsilon M_{\rm box}/3$; the identity of that source may vary
with the point. Good clearance bounds its contribution by $C$, and the
other sources contribute at most
$Ck_{\rm src}h/(\epsilon M_{\rm box})$. Taking the pointwise supremum
therefore gives
\[
 |A(\Lambda)|\le C\left(1+
       \frac{k_{\rm src}h}{\epsilon M_{\rm box}}\right)\le C
\]
for every early atom. This still applies if its first transport occurs
on a later good support, because it tests the lift on $\mathcal C_A$.
At a microscopic edge the same pointwise argument uses
$(1+\operatorname{dist}(v^\flat,x_i))^{-1}\le1$ for the possible near source.
It gives a uniform intrinsic offset for at most two sources at every
volume and for each fixed separated configuration in all sufficiently
large volumes. No local scalar lift is imposed in a defect core; its
intrinsic scalar phases remain real.

For clarity, transport acts on a selected untagged live root by substituting
$z+\Lambda$ for its current field argument, and acts as the identity on an
empty root. Call it an admissible early transport when every induced
offset satisfies the fixed early bounds above; in the separated-source
case this is invoked only when all its atom birth sizes are early.
Write an old atom's factor as
$e^{t[A(Z_<+\tau_A z)+d_A]}$, where $Z_<$ contains its earlier
shell fields and $\tau_A\in[0,1]$ is its current coupling. The factor becomes
\[
 e^{t[A(Z_<+\tau_A z)+d_A+\tau_A A(\Lambda)]}.
\]
The new bounded term stays in the decorated endpoint integrand. It is
not factored into the structural numerical coefficient. The analogous
rule adds the attenuated physical offsets at original microscopic leaves
and a real phase at original charges. A later defect Taylor substitution
$z\mapsto\rho z+t'u$ changes this atom factor to
$e^{t[A(Z_<+\tau_A\rho z)+d_A+\tau_A A(\Lambda)+\tau_A t'A(u)]}$.
The physical tag stays fixed,
the future coupling is multiplied by the real $\rho$, and the new probe
increment is covered by the summable $h/m$ bound above. This proves
closure of the early decoration class under subsequent local operations.

The remaining birth sizes satisfy
$h>c\epsilon M_{\rm box}/[L(1+k_{\rm src})]$. In a finite box, bounded
overlap gives at most $C(M_{\rm box}/h)^2$ atoms of each such size in one
live root, including the fixed reflection multiplicity. There are only
$O_{k_{\rm src},\epsilon}(1)$ such scales. Hence their total number is
at most $C_{k_{\rm src},\epsilon,L}$, independently of the earlier depth.
Good clearance still gives $|A(\Lambda)|\le Ck_{\rm src}$ for each such
event. For these late events only, factor the possibly large deterministic term
$e^{t\tau_A A(\Lambda)}$ into an explicitly recorded late numerical
factor. The variable source remains $tA(Z)$ with the original fixed
contour radius. The product of the late factors is bounded by a constant
depending on the fixed configuration per live root. It does not change
any earlier event weight. Opaque completed coefficients are not
transported, and no singular lift is imposed on a new defect-core tensor.

Finally a fixed-radius polynomial contour produces a bounded linear
Gaussian atom. Within one live root Lemma~\ref{xyh:history} supplies the
overlap hypothesis of Lemma~\ref{xyh:atoms}, so all their source squares
sum to at most $CN_h$. Completing a decorated microscopic tilt produces
sources $R_w^t\Theta a$ or $R_p^t\Theta a$ with bounded coefficients at
each selected edge. Lemma~\ref{xyh:contraction} gives the same square
bound, or these sources may be regarded as edge atoms born at size one.
The deterministic early atom decorations have modulus at most
$e^{C N_h}$ after all fixed-radius contours, while physical phase
decorations have modulus one. These are uniform endpoint bounds for the
specified decoration class, not changes of its structural coefficients.

\begin{lemma}[Estimate for an evaluated history]\label{xyh:endpoint-estimate}
Fix one live Gaussian root with an admissible early decoration record as
defined above. Let $n_F$ be its number of original $F_b$ leaves, and let
$Y_o$ be its real original charge source on its shell stack. Keep all
admissible early decorations inside the endpoint integral. Omit only its
structural numerical coefficient, its opaque completed coefficient
occurrences, its explicitly extracted late numerical factors, and the
phase-probe factors in Equation~\eqref{xyh:phase-probe-cost}. Then the
resulting fully evaluated integral $\mathcal I_h^{(0)}$ obeys
\begin{equation}\label{xyh:endpoint-bound}
 |\mathcal I_h^{(0)}|\le b^{-n_F}\exp\{CN_h-c\|Y_o\|^2\}.
\end{equation}
The same bound holds for normalized fixed-order directional derivatives,
with their bounded number of contour events included in $N_h$. Constants
are uniform over the admissible early decoration class, strict
regularizations, the permitted small parameters, depth, and domain.
The omitted late numerical factors have the separate fixed-configuration
bound proved after Equation~\eqref{xyh:early-source-scales}. Each other
completed root in the forest satisfies its own copy of this estimate.
\end{lemma}
\begin{proof}
The fixed-radius atom contours and the bound $|d_A|\le C_{\rm tag}$ give
$\prod_A|e^{t_A d_A}|\le e^{CN_h}$ within this endpoint integral.
Physical scalar phases have modulus one. The decorated microscopic edge
offsets are uniformly bounded. These facts use exactly the admissible
decoration rules and do not move any early factor into the structural
coefficient.

On a shell acting on current phase charges $k$, the original source is
$\Gamma_r^{1/2}k$ in the $y$ coordinates. It is real because all previous
charge attenuations were real. Write $Y=Y_o+Y_a$; the preceding source
estimates give $\|Y_a\|^2\le CN_h$, and every linear $x$ source has the
same bound. Lemma~\ref{xyh:mixed} reduces the imaginary evaluations at
$F_b$ leaves to the attenuated outputs of $JI$, $I=\Im Y_a$, together
with bounded deterministic offsets. By Lemma~\ref{xyh:moment-bound}, their
values $v_e$ obey
$\sum_e|v_e|^p\le N_h(Cp)^p$, $p\ge2$. Consequently
\begin{align}\label{xyh:imaginary-payment}
 \sum_e |v_e|^2 e^{|v_e|/\sqrt b}
 &\le N_h\sum_{k\ge0}
 \frac{C^{k+2}(k+2)^{k+2}}{k!\,b^{k/2}}
 \le CN_h
\end{align}
for all sufficiently large fixed $b$. The ratio
$(k+2)^{k+2}/k!$ is bounded by an exponential in $k$ times a fixed
polynomial, so the displayed series converges. Bounded offsets add at
most $C$ per leaf. Thus the imaginary growth term in
Equation~\eqref{xyh:primitive-bound} costs $e^{CN_h}$.

Choose $\kappa_1$ strictly between $\kappa_0$ and $1/2$. Small parameter
size and Lemma~\ref{xyh:contraction} give a fixed $\nu>0$ such that
\begin{equation}\label{xyh:gaussian-slack}
 2\kappa_1R_F\{V+C'D^{-1}(C')^t\}R_F^t
   =2\kappa_1R_FM^{-1}R_F^t\preceq(1-\nu)I.
\end{equation}
For a real Gaussian $\xi$ of covariance $\Sigma=R_FVR_F^t$, completing
the square in
$\E e^{\kappa_1\|\xi+R_FC'Y_o\|^2}$ gives determinant
$\det(I-2\kappa_1\Sigma)^{-1/2}$, at most $e^{Cn_F}$.
Its shifted-mean exponent is
$\kappa_1\langle R_FC'Y_o,(I-2\kappa_1\Sigma)^{-1}R_FC'Y_o\rangle$.
Equation~\eqref{xyh:gaussian-slack}, or its Schur complement, bounds this
by $(1-\nu)(Y_o,DY_o)/2$ after decreasing $\nu$ if necessary.
It therefore leaves a fixed part of the charge exponential
$e^{-(Y_o,DY_o)/2}$ unspent.

Means from $Y_a$ and deterministic offsets have square $O(N_h)$ on the
evaluation edges. The difference $\kappa_1-\kappa_0$ absorbs them with
$e^{CN_h}$, using $2|uv|\le\delta u^2+\delta^{-1}v^2$. The same
inequality absorbs cross terms in $-\Re(Y,DY)/2$, and linear real
Gaussian sources, with an arbitrarily small additional fraction of the
retained charge square. Their source squares are $O(N_h)$; Gaussian
integration or a small H\"older enlargement gives $e^{CN_h}$. Since
$D$ is uniformly positive, the remaining suppression is
$e^{-c\|Y_o\|^2}$. Each branch factor contributes the separate
$b^{-1}$ in Equation~\eqref{xyh:primitive-bound}. Polynomial derivatives
use their fixed circles and add only a bounded number of source atoms.
This proves Equation~\eqref{xyh:endpoint-bound} in strict regularization,
with all constants independent of that regularization.
\end{proof}

\subsection{Charge budgets and a composable positive size}

\begin{lemma}[Simultaneous charge payment]\label{xyh:budget}
There are fixed choices of sufficiently large $L,b_0$ and a constant
$c'>0$ such that all phase-probe costs in one generated live root, a specified
bounded final smooth-test cost, and a factor $L^{-4}$ at every small
charged linear step can be included simultaneously in
Equation~\eqref{xyh:endpoint-bound}. Writing $\mathcal I_h$ for the
integral with those phase-probe factors restored, this leaves
\begin{equation}\label{xyh:paid-bound}
 |\mathcal I_h|\le
 C^{N_h}b^{-n_F}L^{-4n_{\rm ch}}
       \exp\{-c'b\textstyle\sum_vq_v^2\}
\end{equation}
apart from its structural numerical coefficient, opaque completed
coefficient occurrences, and explicitly extracted late factors. Here
$q_v$ are the original nonzero integral charge leaves of this root and
$n_{\rm ch}$ counts only its indicated single-input charged steps. The
admissible early decorations remain inside the evaluated integral.
\end{lemma}
\begin{proof}
Fix one selected live root, rather than the sum of its alternatives or
the forest of its separately completed coefficient dependencies.
At a charge-carrying node let $H=h_m(k;X)$ be the incoming test budget.
At a merger budgets add subadditively by test restriction and the
triangle inequality. Reblocking without localization is nonincreasing
by Equation~\eqref{xyh:test-coarsening}. At a neutral Taylor node its
real attenuation $\rho\in[0,1]$ gives
\[
 H_{\rm out}\le\rho(C/L)H,\qquad
 H-H_{\rm out}\ge(1-C/L)H.
\]
Thus a fixed multiple of this decrease pays
Equation~\eqref{xyh:phase-probe-cost}. At termination the remaining budget
is consumed, which also permits bounded final test directions.

For a charged small linear node put $Q=\sum k\ne0$, $q=|Q|$,
$\ell=\log L$, and let $\Gamma$ be its shell. Choose an anchor $v$ in
its support. Lemma~\ref{xyh:kernels}, summing its positive-order binary
bounds, shows that $\Gamma(\cdot,v)-\Gamma(v,v)$ on the local support
and collar is a fixed multiple of a fine test. Consequently
\[
 d:=\Gamma(v,v)\ge c_0\ell,\qquad
 |(k,\Gamma e_v)-Qd|\le C_1H.
\]
Fix $\epsilon=c_0/(2C_1)$ when $C_1>0$. If $H\le\epsilon q\ell$,
positive-semidefinite Cauchy--Schwarz gives
\begin{equation}\label{xyh:energy-alternative}
 (k,\Gamma k)\ge\frac{|(k,\Gamma e_v)|^2}{d}
       \ge\frac{q^2d}{4}\ge c q^2\ell.
\end{equation}
If $C_1=0$ this conclusion holds without a case split. Otherwise, for
$H>\epsilon q\ell$, Equation~\eqref{xyh:test-coarsening} and choices
$C/L\le1/2$, $\ell\ge4/\epsilon$ give
\begin{equation}\label{xyh:potential-alternative}
 H-H_{\rm out}\ge(1-C/L)H-q\ge H/4
                         \ge(\epsilon/4)q\ell.
\end{equation}
At all other nodes the first inequality in
Equation~\eqref{xyh:test-coarsening} is retained; thus no negative
budget decrease is hidden in this dichotomy.

The sum of child budgets minus the parent budget is a nonnegative
potential loss at every node. Summing these losses over this live root,
including the final retained budget, gives at most its original leaf
budget
\begin{equation}\label{xyh:initial-budget}
 H_{\rm leaf}\le C\mathfrak p\sum_v|q_v|.
\end{equation}
Only neutral total charges are multiplied by a Taylor attenuation.
Mergers add totals, so every nonzero sector total remains in
$\mathfrak p\Z$, even though individual entries inside a neutral branch
may cease to be integral multiples. Hence $q\ge\mathfrak p$ at every
recorded charged node. A fixed potential multiplier pays both neutral
probes and $4\ell$ at each node of
Equation~\eqref{xyh:potential-alternative}, once $b$ is large.

The alternative using shell energy requires a separate account. At the
shell of a recorded charged step, its temporary branch was a single
input. Every other live charge group in the eventual combined live root
is either another temporarily separated object at this step or first
joins at a later product ancestor. Lemma~\ref{xyh:history}(a),(d), with
the scalar shell range included in the radius choice, gives zero
cross-covariance between that charged branch and every other such group
on this shell. Thus its scalar energy splits off exactly, with a
nonnegative remaining square. Two selected branches at that shell are
mutually separated as well. Different recursion levels use disjoint
shell coordinates. A later Taylor attenuation changes future fields
only, so it does not alter a payment on an already integrated shell.
Writing $\mathcal E$ for the sum of the selected energies and $E_0$ for
the first-piece square, we have
\begin{equation}\label{xyh:disjoint-reserves}
 E_0=\gamma\mathfrak p^2\sum_vq_v^2,
 \qquad E_0+\mathcal E\le\|Y_o\|^2.
\end{equation}
This is a separation statement about generated single-input histories;
positivity alone would not allow one to split arbitrary charge squares.

Let $c$ be the constant in Equation~\eqref{xyh:endpoint-bound}. Choose a
fixed multiplier $D_0'$ large enough for all potential payments and the
specified final normalized tests. Since $|q_v|\le q_v^2$ for nonzero
integers, increasing $b$ yields
\[
 D_0'H_{\rm leaf}\le(c/4)E_0.
\]
Equations~\eqref{xyh:energy-alternative} and $q\ge\mathfrak p$ similarly
give $4\ell N_E\le(c/4)\mathcal E$, where $N_E$ counts the energy-paid
steps. These fixed allocations use disjoint portions of
Equation~\eqref{xyh:disjoint-reserves}; a fixed positive fraction of
$E_0$ and of the full source square remains. It supplies the microscopic
Gaussian in Equation~\eqref{xyh:paid-bound}. Detached scalar calculations
have their own completed integrals and are absent from this live source
square. Reusing an evaluated scalar multiplies its established bound once
per opaque occurrence; it does not reinsert any charge leaf or shell
reserve. If the occurrences are expanded, apply this proof separately to
each completed root and multiply the resulting bounds. No joint overlap
or charge-square assertion across those roots is used.
\end{proof}

\begin{proposition}[The exact XY activity calculus]\label{xyh:calculus}
Fix the covariance construction and the local partition map above.
Choose the fixed support and collar clearances, small-support threshold,
and $K_0$ first; then choose sufficiently large dyadic $L$, the polymer
weight as in Proposition~\ref{rg:partition-map}, a sufficiently small
parameter neighborhood, and finally sufficiently large $b$. Within
Equation~\eqref{xyh:permitted-parameters} the original XY integrals admit
positive representation sizes with all the following properties. Write
$\mathsf s_X(\mathcal R)$ for the size of a specified representation on
$X$, and $\|\mathcal R\|_A=\sup_X A^{|X|}\mathsf s_X(\mathcal R)$.
\begin{enumerate}[label=\textup{(\roman*)}]
\item An $F_b$ leaf costs $C/b$; a nonzero charge $q$ costs
$C e^{-cbq^2}$, summable over $q\ne0$; the linear microscopic
compensation costs $C|s_\bullet|$, and its nonlinear remainder costs
$C|s_\bullet|^2$. A full-block tensor costs $C|t|$, and a localized
polynomial costs a constant times its coefficient on normalized
$m$-scale directions.
\item Convolving a compatible product multiplies its input sizes, up to
a fixed factor per input. Temporary exclusions, relocation, scalar
normalizer decorations, and separation into a singleton tensor and
remainder are exact. The animal and placement sums obey
Proposition~\ref{rg:partition-map}; hence all multiple-input and
normalizer corrections are quadratic in the input sizes.
\item For a small single neutral linear input of size $u$, localization
has normalized coefficients at most $Cu$ and remainder size
$CL^{-3}u$ in the even bulk after constant and quadratic subtraction,
$CL^{-2}u$ at an even wall after constant subtraction, and $CL^{-1}u$
at a defect after constant subtraction. Each small charged single-input
step has size at most $CL^{-4}u$. Constants in these local gains do not
grow with sufficiently large $L$.
\item Sizes bound terminated evaluations and fixed-order Cauchy
derivatives with bounded smooth normalized backgrounds. In finitely many
normalized source directions the evaluated expressions are jointly
holomorphic on fixed bounded polydiscs, with the corresponding Cauchy
bounds after adding their finitely many probe events. They apply to
copied good pieces, their one-time translations to defects, and the
intrinsic offsets of Equation~\eqref{xyh:offsets}. For fixed separated
sources, the lower bound on $b$ is independent of their number; only the
boundedly many late numerical event costs may depend on the fixed source
configuration. If $T_\Lambda$ is an admissible early transport of an
untagged representation, then for any $A_U\ge A_D>0$, before support relocation
\[
 \mathsf s_X(T_\Lambda\mathcal R)=\mathsf s_X(\mathcal R),\qquad
 A_D^{|X|}\mathsf s_X(T_\Lambda\mathcal R)
 \le (A_D/A_U)^{|X|}\|\mathcal R\|_{A_U}.
\]
The transport leaves every opaque completed coefficient occurrence
unchanged and is not defined on an already tagged live root.
\item Every finite-depth activity, assigned coefficient, and normalized
partition identity is continuous in the permitted parameters, including
at $s=d_*s_{12}^2$ and at $s=0$, as long as its scalar normalizers are
nonzero. The statement holds on the plane at finite support, on finite
reflected grids, and on auxiliary tori with terminal transverse harmonic
directions. The bulk tensor and its local recursion are the same in all
these settings wherever the calculation has the required clearance.
\end{enumerate}
\end{proposition}
\begin{proof}
For a specified representation write each selected term in the form
\[
 c_h(\alpha)\,\ell_h(\alpha)
       \left(\prod_{\nu=1}^{m_h}a_{h,\nu}\right)
       \mathcal I_h(z;\mathbf d_h,\alpha).
\]
Here $h$ is its unique live root, or the empty root with integral one.
The variables $\alpha$ range over the compact Taylor and contour
parameter spaces with their positive product measure. The structural
coefficient $c_h$ contains the displayed parameter powers, finite
averaging coefficients, Fourier/Taylor direction weights, and contour
kernel factors. The factor $\ell_h$ is only the explicitly extracted
late numerical product defined above, and equals one when there are no
such events. The $a_{h,\nu}$ are opaque occurrences of previously
completed coefficients. The admissible early record $\mathbf d_h$ and
the generated phase-probe factors paid by Lemma~\ref{xyh:budget} stay
inside $\mathcal I_h$ and are not part of $c_h$ or $\ell_h$.

Give the live root the weight
\begin{equation}\label{xyh:representation-weight}
 W(h)=C_*^{N_h}\,b^{-n_F}
 \exp\{-c_*b\textstyle\sum_vq_v^2\}\,L^{-4n_{\rm ch}},
\end{equation}
with $W=1$ for the empty root. All counts and charges in this display
belong only to $h$. If $a$ has a chosen completed representation, let
$\mathsf m(a)$ be the bound supplied by this same rule at its earlier
evaluation, on the source domain where the size will be used. A size at
one fixed source point may instead use the exact modulus $|a|$ there.
For a uniform estimate or a later Cauchy formula on a source polydisc,
$\mathsf m(a)$ must be a completed size uniform on that polydisc, or
another proved bound for the supremum of $|a|$ there; a modulus at one
point does not suffice. Equivalently one may differentiate $a$ in its
separately completed root and use that derivative's completed size.
These choices also apply when completed numbers are combined before
forming a new primitive. In particular each source-independent component
of the full unmarked tensor may use
$\mathsf m(t_{ab,j})=|t_{ab,j}|\le |t_j|_\infty$ after its evaluated
increments have been summed, preserving their cancellation. This choice
asserts no derivative bound in the real reference parameters; only their
continuity is used in (v). The acyclic coefficient dependence makes this
a finite-depth recursive definition. Set
\[
 \mathsf s_X(\mathcal R)=\sum_h\int
     |c_h(\alpha)|\,|\ell_h(\alpha)|
     \left(\prod_{\nu=1}^{m_h}\mathsf m(a_{h,\nu})\right)
     W(h)\,d\alpha.
\]
Choose $C_*$ sufficiently large and $c_*>0$ sufficiently small from the
bounds in Lemmas~\ref{xyh:endpoint-estimate} and~\ref{xyh:budget}, which
are uniform over the admissible early records and sufficiently large
$L$. Induction over coefficient evaluations then gives
$|a|\le\mathsf m(a)$ and makes this size a positive majorant of every
endpoint forest. No arbitrary numerical factor can be hidden in the
endpoint integral because its admissible records have the specified
generation rules and bounds.

The primitive bounds in (i) follow directly: polynomial
contours contribute a fixed number of atoms, and
Equation~\eqref{xyh:quadratic-remainder} has two explicit powers of
$s_\bullet$ and only a bounded tilt strength in its exponential.
A full tensor has normalized vertex averages: each selected summand
has two normalized atoms and coefficient $O(|t|)$.

When compatible live roots are concatenated, their live leaf lists and
actual event lists are disjoint. The common-stack construction joins them
into one live root, while their structural and late numerical factors
multiply and their lists of opaque coefficient occurrences concatenate.
At the final endpoint apply Equation~\eqref{xyh:paid-bound} once to the
concatenated live root, whose geometric properties were proved in
Lemma~\ref{xyh:history}. Thus no fresh $C^{N_h}$ is charged each time an
old root is convolved. Each opaque numerical coefficient is already
bounded at its own endpoint. In particular an inverse-normalizer term
$a_B^m$ has size $\mathsf m(a_B)^m$, and its absolute series is at most
$\mathsf m(a_B)/(1-\mathsf m(a_B))$ when $\mathsf m(a_B)<1$. If these
occurrences are expanded, they are $m$ separately completed roots and
their endpoint weights multiply; no common atom-overlap estimate is
applied to them.
This proves the product rule. The purely spatial enumeration in
Proposition~\ref{rg:partition-map} uses precisely these multiplicative
sizes: nonlinear inputs carry at least two factors, large supports use
the polymer weight gain, and the inverse scalar decorations use their
absolutely convergent small scalar series. It does not require a
pointwise norm at arbitrary remaining fields. This proves (ii).

For (iii) Taylor's integral formula of order $a$ contributes $a$ direction
arguments and a fixed number of contour events. Apply the exact
$CL^{-1}$ directional decomposition of Lemma~\ref{xyh:fourier} to each
argument before using Cauchy's formula. This gives $CL^{-a}$ in the
positive coefficient sum. The evenness assumption makes the linear
term vanish when $a=2$ or $3$. In the bulk, the quadratic Taylor term
has two arguments; replacing one by its affine approximation costs
$CL^{-2}$ and the other $CL^{-1}$, so the two-factor difference has
$CL^{-3}$. The affine polynomial itself is the assigned tensor
prescription of Equation~\eqref{xyh:tensor}. Its normalized coefficient
is a terminated history with bounded probes, hence costs $Cu$ by
Lemma~\ref{xyh:budget}. Scalar extraction is the same statement with
zero Taylor arguments. The charged gain is exactly the recorded factor
in Equation~\eqref{xyh:representation-weight}. Multi-input terms receive
no such Taylor or charged gain, consistently with both the grammar and
the quadratic counting rule.

For (iv), apply $T_\Lambda$ only to an untagged live root in the displayed
representation, and as the identity to an empty root, using the same
source domain and completed coefficient bounds. The transport rule changes
$\mathbf d_h$ inside $\mathcal I_h$ and leaves $c_h,\ell_h$, every $\mathsf m(a_{h,\nu})$,
and the event list defining $W(h)$ unchanged. Its own-birth bound proves
that the new record is admissible. The equality of representation sizes
in (iv) is therefore termwise in their defining nonnegative sums, before
any support relocation; multiplying by $A_D^{|X|}$ gives the displayed
comparison with the $A_U$ norm. In particular no factor per old event is
charged to a support weight. The endpoint bound for the transported
integrand is available because Lemma~\ref{xyh:endpoint-estimate} was
proved uniformly over these records. This is a comparison of the same
chosen representations, not merely a comparison of their endpoint values.
At finite strict regularization the factor-by-factor substitution is an
exact representation of the translated activity. The uniform majorants
used below pass this identity to the limiting evaluated expressions on
the bounded background domains used here.

The derivative claims in (iv) follow from the derivative version of
Lemma~\ref{xyh:endpoint-estimate} and the final-test allowance in
Lemma~\ref{xyh:budget}. A Cauchy derivative in a remaining field
direction acts on $\mathcal I_h$ with every opaque scalar occurrence
fixed. A derivative in a source parameter on which an opaque scalar
depends instead uses its uniform polydisc bound in the preceding
definition, or differentiates that scalar within its separately completed
root. Thus cutting a root never discards a source-parameter derivative.
The summable later-probe calculation proves that
a tagged root stays in the admissible class under further defect
operations. At late birth sizes only the separate factor $\ell_h$ changes;
its bounded number of events gives a fixed-configuration factor per live
root. It does not change an earlier event weight or an opaque completed
coefficient size. Thus the inverse-temperature threshold remains
independent of the fixed correlation order.

Finally fix a finite calculation depth and a finite geometry. All its
support choices and discrete Fourier direction decompositions are
finite; compact Taylor and contour parameters can be fixed independently
of a small parameter variation. Only the microscopic integral charge
indices remain unbounded. Before removing regularizations, all
identities are ordinary integrated identities of entire functions with
strict real-density integrability: the $x$ growth is below its Gaussian
threshold by Equation~\eqref{xyh:primitive-bound}, the real $y$ precision
is positive with strict compensation, and the mixed terms are imaginary
on real coordinates except for controlled lower-degree sources.
Lemma~\ref{xyh:mixed} proves the integrated formulas and
Lemmas~\ref{xyh:endpoint-estimate}--\ref{xyh:budget} give bounds uniform
as compensation becomes full. At fixed geometry the completion before
the Cameron estimate is also continuous: its real covariance $V$ may
become singular, but positive Gaussian measures converge to the measure
on its supported space, and its complex shifts are bounded for fixed
contour choices. The retained real-charge Gaussian square dominates the
remaining fixed-background costs, which grow only exponentially linearly
in the charge magnitudes. Thus dominated convergence removes charge
truncation and strict compensation at this fixed depth. The same
majorants hold uniformly on each smaller source polydisc, so integrating
Cauchy formulas proves holomorphy and convergence of source derivatives.
It also passes real parameter continuity to the limit. Scalar division is
continuous whenever its normalizer stays away from zero, as ensured in
the small-activity map. Plane calculations at finite depth use an
auxiliary large torus whose preterminal covariance entries agree on
all required supports, so the same reasoning applies there. This proves
(v) and completes the limiting justification promised in
Proposition~\ref{xyh:compensation}.
\end{proof}
